\documentclass[a4paper,11pt,openany]{report}
\usepackage{../preamble/agenticboost}

\title{%
  \textsc{AgenticCareerBoost}\\[0.4em]
  \Large Agentic System Guide\\[0.3em]
  \large Human-facing orientation manual%
}
\author{Dídac Llorens \\ \small \href{https://github.com/DidacLL}{github.com/DidacLL}}
\date{April 2026 --- v1.0}

\begin{document}
\maketitle
\tableofcontents
\newpage

% ============================================================================
\chapter{What This Project Is}
% ============================================================================

\takeaway{AgenticCareerBoost is a public engineering campaign built as a
path-based multiagent operating system inside a Git repository. The repository
is not just where the project is documented; it is where the system runs.}

The project exists to turn work into inspectable proof. Instead of describing a
technical profile with adjectives, it records decisions, workflows, state, and
formal outputs directly in versioned files. Humans can read the repository as a
project manual. Agents can read it as an execution environment.

This guide is the human-readable entrypoint. The deeper technical evidence is in
the bootstrap case study:
\pathref{content/reports/build/s000-agentic-os-bootstrap.pdf}.

\section{What it is not}

\begin{itemize}
  \item Not a startup product.
  \item Not a generic ``AI wrapper'' demo.
  \item Not a private automation setup hidden outside the repo.
  \item Not a replacement for human judgment.
\end{itemize}

\pitfall{Do not mistake the repository structure for bureaucracy. The explicit
files exist so humans and models can inspect the same truth with minimal
ambiguity.}


% ============================================================================
\chapter{Mission, Purpose, and Intended Audience}
% ============================================================================

\takeaway{The mission is to rebuild a public technical profile through visible,
agentic engineering work and make that work readable by recruiters, peers, and
future collaborators.}

\section{Mission and purpose}

The canonical mission lives in \pathref{docs/core/mission.md}. In practical
terms, the system has three public jobs:

\begin{itemize}
  \item convert scattered capability into clear engineering artifacts,
  \item keep repository, site, and social outputs coherent,
  \item prove the process, not only the result.
\end{itemize}

\section{Intended audience}

\begin{description}[style=nextline]
  \item[Curious readers] need a short explanation of what the repository is and
  why it exists.
  \item[Recruiters and hiring managers] need fast, credible evidence of
  systems thinking, execution discipline, and documentation depth.
  \item[Engineering peers] need the actual mechanism: roles, workflows,
  review loops, build paths, and state handling.
  \item[Operators and contributors] need to know how to request work, where to
  read next, and which files are authoritative.
\end{description}

\section{How this guide fits the rest of the docs}

\begin{itemize}
  \item This guide explains the system in plain operational language.
  \item The README is the shorter repo entrypoint.
  \item The site is the public mirror.
  \item The S-000 case study is the detailed technical record.
\end{itemize}


% ============================================================================
\chapter{The System Mental Model}
% ============================================================================

\takeaway{The system has four layers: stable truth, workflow contracts, volatile
state, and published outputs. Agents are routed by paths, not by a giant hidden
prompt.}

\begin{figure}[htbp]
\centering
\begin{tikzpicture}[node distance=1.5cm and 2.1cm]
  \node[usernode] (human) {Human / agent};
  \node[agentnode, below=of human] (entry) {AGENTS.md};

  \node[filenode, below left=1.6cm and 2.7cm of entry] (core) {docs/core/};
  \node[agentnode, below=of entry] (workflows) {docs/workflows/};
  \node[agentnode, below right=1.6cm and 2.7cm of entry] (roles) {docs/agents/};

  \node[statenode, below left=1.8cm and 1.5cm of workflows] (state) {state/};
  \node[filenode, below=of workflows] (reports) {content/reports/};
  \node[filenode, below right=1.8cm and 1.5cm of workflows] (site) {site/ + data/};

  \draw[flowedge] (human) -- (entry);
  \draw[dataedge] (entry) -- node[above left,font=\scriptsize]{truth} (core);
  \draw[flowedge] (entry) -- node[left,font=\scriptsize]{choose workflow} (workflows);
  \draw[flowedge] (entry) -- node[above right,font=\scriptsize]{choose role} (roles);
  \draw[dataedge] (workflows) -- (state);
  \draw[dataedge] (roles) -- (state);
  \draw[flowedge] (workflows) -- (reports);
  \draw[flowedge] (reports) -- (site);
\end{tikzpicture}
\caption{System mental model: entrypoint, contracts, state, and outputs.}
\label{fig:mental-model}
\end{figure}

\section{Stable truth}

Stable truth lives under \pathref{docs/core/}. These files describe the mission,
constraints, brand, and the truth hierarchy. They change rarely and win when a
volatile file disagrees.

\section{Workflow contracts}

Workflow files define the allowed action modes:
\workflow{chat}, \workflow{operate}, \workflow{review}, \workflow{hotfix},
\workflow{plan}, \workflow{sprint}, and \workflow{system-review}.

\section{Volatile state}

Volatile state lives under \pathref{state/}. It holds the current status,
backlog, active sprint, logs, and family memory. It is important, but it is not
canonical truth when it conflicts with core rules.

\section{Published outputs}

Outputs live under \pathref{content/}, \pathref{site/starter/}, and
\pathref{data/}. The formal PDFs are part of that surface: they are meant to be
read, cited, and inspected.


% ============================================================================
\chapter{How To Use The System As A Human}
% ============================================================================

\takeaway{Humans use the repository by choosing the right workflow, not by
reading every file. Start from the question you are trying to answer or the
change you are trying to make.}

\section{Use it as a curious reader}

If you want to understand the project quickly:

\begin{enumerate}
  \item Read the README.
  \item Read this guide.
  \item Open the S-000 case study only if you want deeper evidence.
\end{enumerate}

\section{Use it as an operator or collaborator}

If you want to get work done:

\begin{enumerate}
  \item Start at \pathref{AGENTS.md}.
  \item Choose the workflow that matches the request.
  \item Read only the referenced contracts and current state.
  \item Keep changes inside declared scope and writes.
\end{enumerate}

\section{Use it as an agent}

Agents should not improvise the filesystem. The correct path is:
\pathref{AGENTS.md} $\rightarrow$ workflow $\rightarrow$ role or AutoAgent
$\rightarrow$ state and outputs.

\begin{figure}[H]
\centering
\begin{tikzpicture}[node distance=1.4cm and 1.8cm]
  \node[usernode] (request) {Request};
  \node[stepnode, below=of request] (entry) {Read AGENTS.md};
  \node[stepnode, below=of entry] (choose) {Choose workflow};
  \node[stepnode, below left=1.7cm and 2.1cm of choose] (single) {One specialist};
  \node[stepnode, below right=1.7cm and 2.1cm of choose] (sprint) {Multi-step sprint};
  \node[statenode, below=of single] (output1) {Scoped output};
  \node[statenode, below=of sprint] (output2) {Reviewed artifacts};

  \draw[flowedge] (request) -- (entry);
  \draw[flowedge] (entry) -- (choose);
  \draw[flowedge] (choose) -- node[above left,font=\scriptsize]{operate / hotfix / review} (single);
  \draw[flowedge] (choose) -- node[above right,font=\scriptsize]{plan / sprint} (sprint);
  \draw[dataedge] (single) -- (output1);
  \draw[dataedge] (sprint) -- (output2);
\end{tikzpicture}
\caption{Human usage flow: choose the workflow before choosing the files.}
\label{fig:human-usage-flow}
\end{figure}


% ============================================================================
\chapter{Actions And When To Use Them}
% ============================================================================

\takeaway{The action surface is intentionally small. Each workflow exists to
solve a different coordination problem.}

\section{Action quick reference}

\begin{longtable}{p{2.2cm} p{4.2cm} p{5.6cm}}
\toprule
\textbf{Action} & \textbf{Use it when} & \textbf{Typical result} \\
\midrule
\endhead
\workflow{chat} & You need discussion only & Answers, constraints, optional summary \\
\workflow{operate} & One bounded specialist should do one bounded job & One scoped artifact or correction \\
\workflow{review} & You need maintenance checks and consistency fixes & Scoped sync, lint, compile verification \\
\workflow{hotfix} & One small fix is clearly isolated & One commit and minimal backlog note \\
\workflow{plan} & A larger change needs a contract first & Active sprint definition \\
\workflow{sprint} & Planned multi-step work needs role coordination & Reviewed multi-artifact delivery \\
\workflow{system-review} & The architecture or rule system itself looks wrong & Issue report and proposed structural changes \\
\bottomrule
\end{longtable}

\section{Action selection flow}

\hfigure{
	\begin{tikzpicture}
	  \node[usernode] (start) {Need};
	  \node[stepnode, below=of start] (question1) {Only discussion?};
	  \node[stepnode, below left=1.5cm and 2.0cm of question1] (chat) {\workflow{chat}};
	  \node[stepnode, below right=1.5cm and 2.0cm of question1] (question2) {One bounded task?};
	  \node[stepnode, below left=1.5cm and 2.0cm of question2] (operate) {\workflow{operate}};
	  \node[stepnode, below right=1.5cm and 2.0cm of question2] (question3) {Needs planning / multiple roles?};
	  \node[stepnode, below left=1.5cm and 2.0cm of question3] (review) {\workflow{review} or \workflow{hotfix}};
	  \node[stepnode, below right=1.5cm and 2.0cm of question3] (plan) {\workflow{plan} $\rightarrow$ \workflow{sprint}};

	  \draw[flowedge] (start) -- (question1);
	  \draw[flowedge] (question1) -- node[above left,font=\scriptsize]{yes} (chat);
	  \draw[flowedge] (question1) -- node[above right,font=\scriptsize]{no} (question2);
	  \draw[flowedge] (question2) -- node[above left,font=\scriptsize]{yes} (operate);
	  \draw[flowedge] (question2) -- node[above right,font=\scriptsize]{no} (question3);
	  \draw[flowedge] (question3) -- node[above left,font=\scriptsize]{maintenance or tiny fix} (review);
	  \draw[flowedge] (question3) -- node[above right,font=\scriptsize]{planned multi-step work} (plan);
	\end{tikzpicture}
}{Choose the workflow by scope and coordination needs.}{action-flow}

% ============================================================================
\chapter{Roles, Specialties, AutoAgents, And Memory}
% ============================================================================

\takeaway{Top-level roles stay small; specialties narrow execution; AutoAgents
handle recurring routines; family memory stores durable heuristics, not truth.}

\section{Human-facing roles}

\begin{itemize}
  \item \role{Orchestrator} routes work and writes contracts.
  \item \role{Developer} implements bounded technical changes.
  \item \role{PairCheck} reviews non-trivial outputs independently.
  \item \role{CI/CD} integrates and publishes approved work.
  \item \role{Documentation} turns changes into readable engineering artifacts.
  \item \role{CommunityManager} adapts proof into public-facing messaging.
\end{itemize}

\section{Specialties}

Roles should not execute as generic experts. They run in specialty modes such
as \texttt{Developer/latex}, \texttt{Documentation/guide}, or
\texttt{CommunityManager/linkedin}. Specialty keeps the read scope small and
the expected output explicit.

\section{AutoAgents}

The fixed routines currently live in \pathref{docs/agents/autoagents.md}. They
cover consistency, Markdown quality, LaTeX compilation, and social planning.

\section{Family memory}

Family memory lives in \pathref{state/memory/}. It stores durable heuristics
such as recurring build issues, accepted phrasing defaults, or source-quality
notes. It must never replace canonical truth in \pathref{docs/core/}.


% ============================================================================
\chapter{Repository Layout And Where To Start}
% ============================================================================

\takeaway{The repository is organized so that each question leads to one or two
obvious next files.}

\section{Condensed layout}

\begin{tabularx}{\linewidth}{l X}
\toprule
\textbf{Path} & \textbf{Why you would open it} \\
\midrule
\pathref{AGENTS.md} & Start here if you are routing work \\
\pathref{docs/core/} & Understand what is stable and non-negotiable \\
\pathref{docs/workflows/} & Choose the correct action mode \\
\pathref{docs/agents/} & Understand who owns the work \\
\pathref{state/} & Read current status, backlog, and recent activity \\
\pathref{content/reports/} & Read published formal documentation \\
\pathref{site/starter/} & Inspect the public mirror \\
\bottomrule
\end{tabularx}

\section{Where a new reader should start}

\begin{enumerate}
  \item README for the shortest orientation.
  \item This guide for the operating model.
  \item S-000 for technical depth.
  \item Specific workflows and roles only when you need to act.
\end{enumerate}


% ============================================================================
\chapter{Typical Flows}
% ============================================================================

\takeaway{Most real usage falls into a few repeatable flows: read, operate,
review, or plan.}

\section{Human orientation flow}

README $\rightarrow$ this guide $\rightarrow$ S-000 case study.

\section{Bounded maintenance flow}

\workflow{review} $\rightarrow$ \texttt{ContentSync} first
$\rightarrow$ Markdown or LaTeX specialist if needed
$\rightarrow$ scoped fixes only.

\section{Change delivery flow}

\workflow{operate} for one bounded specialist task, or
\workflow{plan} $\rightarrow$ \workflow{sprint} when the work spans multiple
concerns and needs pair-check.

\section{Documentation flow}

Implementation change $\rightarrow$ documentation artifact
$\rightarrow$ published PDF or repo/site update
$\rightarrow$ optional public adaptation by \role{CommunityManager}.

\begin{figure}[htbp]
\centering
\begin{tikzpicture}[node distance=1.4cm and 1.8cm]
  \node[filenode] (readme) {README.md};
  \node[filenode, right=2.8cm of readme] (guide) {agentic-system-guide.pdf};
  \node[filenode, right=2.8cm of guide] (case) {S-000 case study};
  \node[filenode, below=of guide] (site) {site/starter/};
  \node[filenode, below=of case] (build) {content/reports/build/};

  \draw[flowedge] (readme) -- node[above,font=\scriptsize]{quick entry} (guide);
  \draw[flowedge] (guide) -- node[above,font=\scriptsize]{technical depth} (case);
  \draw[dataedge] (guide) -- node[right,font=\scriptsize]{mirrored summary} (site);
  \draw[dataedge] (case) -- node[right,font=\scriptsize]{published PDFs} (build);
\end{tikzpicture}
\caption{Documentation map: quick entrypoint, human guide, technical case study, and published outputs.}
\label{fig:documentation-map}
\end{figure}


% ============================================================================
\chapter{Reading Path To Deeper Technical Material}
% ============================================================================

\takeaway{Once the high-level model is clear, the deepest material is organized
by question, not by chronology.}

\begin{description}[style=nextline]
  \item[What is always true?] Read \pathref{docs/core/}.
  \item[What action should be used?] Read \pathref{docs/workflows/}.
  \item[Who owns each kind of work?] Read \pathref{docs/agents/}.
  \item[What changed recently?] Read \pathref{state/current.md} and the
  backlog.
  \item[How was the bootstrap system implemented?] Read the S-000 case study.
  \item[How are formal PDFs built?] Read \pathref{content/reports/tex/README.md}
  and the local build scripts.
\end{description}

\pitfall{Do not start from logs unless you are debugging a specific failure.
Logs are volatile evidence, not the best orientation path.}


% ============================================================================
\appendix
\chapter{Quick Reference}
% ============================================================================

\takeaway{These are the highest-value paths for humans using the repository.}

\begin{longtable}{p{5cm} p{8cm}}
\toprule
\textbf{Path} & \textbf{Use it for} \\
\midrule
\endhead
\pathref{README.md} & quickest human overview \\
\pathref{AGENTS.md} & routing entrypoint for agents and operators \\
\pathref{docs/core/mission.md} & purpose and success criteria \\
\pathref{docs/workflows/review.md} & maintenance review chain \\
\pathref{docs/workflows/operate.md} & one bounded specialist task \\
\pathref{docs/agents/autoagents.md} & recurring routines and triggers \\
\pathref{state/current.md} & current status snapshot \\
\pathref{content/reports/build/agentic-system-guide.pdf} & published human-facing guide \\
\pathref{content/reports/build/s000-agentic-os-bootstrap.pdf} & published technical case study \\
\bottomrule
\end{longtable}

\end{document}
